\subsection{连续的实值函数}

除了取值于任意拓扑空间的连续函数的一般性质（第一章，§ 2）之外，连续实值函数还具有下面两个基本性质：

定理 1（Weierstrass）. 设 \(f\) 是定义在非空拟紧空间 \(\mathbf{X}\) 上的连续实值函数。则至少存在一点 \(a \in \mathbf{X}\) 使得 \(f(a) = \sup_{x \in \mathbf{X}} (f(x))\) ，且至少存在一点 \(b \in \mathbf{X}\) 使得 \(f(b) = \inf_{x \in \mathbf{X}} f(x)\) 。因为 \(f(\mathbf{X})\) 是紧的（第一章，§ 9，第 4 节，定理 2），从而在 \(\overline{\mathbf{R}}\) 中是闭的；于是 \(f(\mathbf{X})\) 包含它的界。

这个定理常叙述成如下形式：非空拟紧空间上的连续实值函数达到它的界。

推论. 若定义在非空拟紧空间 X 上的实值函数在 X 上连续且有限，则它在 X 中有界。

定理 2（Bolzano）. 设 \(f\) 是定义在连通空间 \(\mathbf{X}\) 上的连续实值函数。若 \(a\) 与 \(b\) 是 \(\mathbf{X}\) 的任意两点，且 \(\alpha\) 是属于以 \(f(a)\) 与 \(f(b)\) 为端点的闭区间的实数，则至少存在一点 \(x \in \mathbf{X}\) 使得 \(f(x) = \alpha\) 。

因为 \(f(\mathbf{X})\) 是连通的（第一章，§ 11，第 2 节，命题 4），从而是 \(\overline{\mathbf{R}}\) 中的一个区间（ \(\S 4\) ，第 2 节，命题 5）；于是它包含以 \(f(a)\) 与 \(f(b)\) 为端点的闭区间。

这个定理常表述成如下形式：连通空间上的连续实值函数从一个值变到另一个值时不可能不经过一切中间值。

这个性质决不是连续函数的特征；存在定义在连通空间上且在每一点都不连续的函数具有此性质（习题 2）。

